\chapter{完成情况及性能参数}
阐述最终实现的成果，图文结合，以实物照片为主。
\section{整体介绍}
提供整个系统实物的正面照片和斜 \(45^\circ\) 全局照片，说明系统完成程度及运行状态。
\section{工程成果}
按硬件实物、软件界面等类别展示设计结果。
\subsection{机械成果}
给出机械结构实物照片并说明装配情况；若不涉及机械设计，明确说明。
\subsection{电路成果}
给出电路实物照片，说明电路板、FPGA 平台、外设和连接情况。
\subsection{软件成果}
给出软件界面照片或运行截图，并说明主要操作、输入输出及运行结果。
\section{特性成果}
逐个展示功能、性能参数等量化指标，可加入重要仪器测试或现场照片。记录测试环境、测试方法、输入条件、实测结果及误差分析，使结果能够复核。

对于 FPGA 设计，可根据作品情况补充时钟频率、吞吐率、处理延迟、逻辑与存储资源占用、功耗等数据，并注明工具和约束条件。
